\section{Omnibus VPC Significance Testing}
\label{sec:global_lrt}

\subsection{Omnibus VPC hypothesis and its covariance representation}
\label{subsec:global_hypothesis}
The variance partition coefficient $\rho(x;\theta)$ describes the
proportion of marginal response variation attributable to between-cluster
heterogeneity at covariate value $x$. We first ask whether this
contribution vanishes throughout the covariate domain $\mathcal X$.
The \emph{omnibus VPC significance test} concerns
\begin{equation}\label{eq:sec5-vpc-null}
 \begin{aligned}
 H_0^{\mathrm{omni}}:&\quad
 \rho(x;\theta)=0\quad(\forall x\in\mathcal X),\\
 H_1^{\mathrm{omni}}:&\quad
 \rho(x;\theta)>0\quad(\exists x\in\mathcal X).
 \end{aligned}
\end{equation}
This is distinct from pointwise VPC significance, which concerns
$\rho(x_0;\theta)=0$ at one prespecified covariate value, and from
VPC equality, which concerns $\rho(x_i;\theta)=\rho(x_j;\theta)$.

Write $z(x)=(1,x)^{\mathsf T}$ and let
$\Gamma\succeq0$ denote the random-intercept-and-slope covariance
matrix. The variance of the random-effects contribution to the linear
predictor is $v(x)=z(x)^{\mathsf T}\Gamma z(x)$.
Throughout this section we use the zero-VPC property established for
the VPC definition in this manuscript: at the covariate values under
consideration, $\rho(x;\theta)=0$ if and only if $v(x)=0$.
In particular, the observation-level variance contribution is assumed
strictly positive so that the VPC is well defined.

\begin{proposition}\label{prop:sec5-omnibus-equivalence}
Suppose $\Gamma\succeq0$, the zero-VPC property holds for all
$x\in\mathcal X$, and $\mathcal X$ contains at least two distinct
covariate values. Then
\[
 \rho(x;\theta)=0\ \text{for every }x\in\mathcal X
 \quad\Longleftrightarrow\quad \Gamma=0.
\]
In particular, the equivalence holds for the binary domain
$\mathcal X=\{0,1\}$.
\end{proposition}

\begin{proof}
For a positive-semidefinite matrix, $z^{\mathsf T}\Gamma z=0$
if and only if $\Gamma z=0$, as follows from a positive-semidefinite
square root of $\Gamma$. If the VPC vanishes at two distinct values
$x_a\ne x_b$, then $\Gamma z(x_a)=\Gamma z(x_b)=0$.
The vectors $z(x_a)$ and $z(x_b)$ are linearly independent, hence
span $\mathbb R^2$, so $\Gamma=0$. The converse is immediate from
the zero-VPC property.
\end{proof}

Accordingly, under the stated design and VPC conditions, the scientific
hypothesis in \eqref{eq:sec5-vpc-null} has the equivalent covariance
representation
\begin{equation}\label{eq:sec5-null}
 \begin{aligned}
 H_0^{\mathrm{omni}}:&\quad\Gamma=0,\\
 H_1^{\mathrm{omni}}:&\quad\Gamma\succeq0,\quad\Gamma\ne0.
 \end{aligned}
\end{equation}
Thus the covariance restriction is a computational representation of an
omnibus VPC hypothesis, not a change of inferential target.
If $\mathcal X$ contains only one distinct value, this equivalence
need not hold: a nonzero rank-one covariance matrix may annihilate
that design vector.

Write $\theta=(\lambda^{\mathsf T},\gamma^{\mathsf T})^{\mathsf T}$,
where $\lambda\in\Lambda\subset\mathbb R^{p_\lambda}$ collects the
unrestricted nuisance coordinates and
$\gamma=(\gamma_{11},\gamma_{12},\gamma_{22})^{\mathsf T}$
parameterizes $\Gamma$, as in Section~\ref{sec:geometry_asymptotics}.
Under the omnibus null $\gamma=0$; under the alternative $\gamma$
belongs to the positive-semidefinite cone $K$.

\subsection{Likelihood-ratio statistic}
\label{subsec:global_lrt_statistic}
Let $\widehat\theta_n=(\widehat\lambda_n,\widehat\gamma_n)$
maximize the marginal likelihood over $\Lambda\times K$, and let
$\widetilde\lambda_n$ maximize it subject to $\gamma=0$.
The likelihood-ratio statistic for the omnibus VPC hypothesis is
\begin{equation}\label{eq:sec5-lrt}
 \begin{aligned}
 \mathrm{LRT}_{\mathrm{omni},n}
 &=2\bigl\{\ell_n(\widehat\theta_n)\\
 &\qquad-\ell_n(\widetilde\lambda_n,0)\bigr\}.
 \end{aligned}
\end{equation}
The statistic is nonnegative for exact global maximizers because the
restricted parameter space is contained in the full parameter space.
Numerically negative values indicate optimization error and should
not be interpreted as realizations from the null law. In particular,
truncation of materially negative numerical values to zero does not
establish that the computed likelihood ratio approximates the exact
likelihood ratio. The asymptotic results below require both fitted
likelihood values to approximate their respective global maxima with
error $o_p(1)$ on the log-likelihood scale.

\subsection{Efficient information under the omnibus VPC null}
\label{subsec:global_efficient_information}
Let $I_0$ be the normalized full Fisher information at
$\theta_0=(\lambda_0,0)$, partitioned as in
Section~\ref{subsec:fisher_projection}. Define
\begin{equation}\label{eq:sec5-J}
 \begin{aligned}
 J_0={}&I_{\gamma\gamma,0}\\
 &-I_{\gamma\lambda,0}
 I_{\lambda\lambda,0}^{-1}I_{\lambda\gamma,0}.
 \end{aligned}
\end{equation}
The efficient covariance score has limiting distribution
$Z_{\gamma\cdot\lambda}\sim N_3(0,J_0)$.
The efficient information $J_0$, rather than the unadjusted
covariance block, accounts for estimation of the unrestricted nuisance
coordinates. Section~\ref{sec:information} develops its evaluation
under the restricted null.


\subsection{Gaussian projection and chi-bar-square limit}
\label{subsec:global_chibar_limit}

Likelihood-ratio statistics under inequality and
convex-cone constraints generally need not follow
the ordinary chi-square limits of unconstrained
likelihood theory \citep{self1987asymptotic,shapiro1985}.
For convex cones, the limiting distribution can
be expressed through Gaussian projections and
associated chi-bar-square mixtures
\citep{amelunxen2014}.

Let $J_0^{1/2}$ be the symmetric positive-definite
square root and put
\begin{equation}
\mathcal C_0=J_0^{1/2}K.
\label{eq:sec5-transformed-cone}
\end{equation}
For a closed convex cone $C$, denote Euclidean
projection onto $C$ by $\Pi_C$.


\begin{theorem}\label{thm:sec5-global-lrt}
Suppose $H_0$ holds at $\theta_0=(\lambda_0^{\mathsf T},0^{\mathsf T})^{\mathsf T}$,
with $I_0\succ0$, and the uniform local quadratic expansion,
feasible-set convergence, and joint full/restricted localization
and approximate-maximization conditions of
Section~\ref{sec:geometry_asymptotics} and
Appendix~\ref{app:asymptotics} hold. Then
\begin{equation}\label{eq:sec5-lrt-limit}
 \mathrm{LRT}_{\mathrm{omni},n}\xrightarrow{d}
 \|\Pi_{\mathcal C_0}(G)\|^2,
 \qquad G\sim N_3(0,I_3).
\end{equation}
Equivalently, the limiting law is the chi-bar-square mixture
\begin{equation}\label{eq:sec5-chibar}
 \sum_{j=0}^{3}w_j\chi_j^2,
 \qquad w_j\geq0,\quad\sum_{j=0}^{3}w_j=1,
\end{equation}
where $\chi_0^2$ is a point mass at zero and
$w_j=v_j(\mathcal C_0)$ are the conic intrinsic volumes of
$\mathcal C_0$.
\end{theorem}

The result follows from joint convergence of the full and restricted
maximized local likelihoods, both driven by the same Gaussian score.
Profiling the common quadratic likelihood over $\lambda$ and
comparing its maximum over $K$ with its maximum at the vertex yields
the Gaussian projection representation. The detailed derivation and
chi-bar-square identification are given in
Appendix~\ref{app:global_lrt}. In particular, an ordinary
$\chi_3^2$ reference distribution does not generally give the
correct null calibration.


\subsection{Exact mixture weights for the
$2\times2$ PSD cone}
\label{subsec:global_exact_weights}

The chi-bar-square weights are the conic intrinsic
volumes of the information-transformed covariance
cone. Their geometric interpretation and the
conic Gauss--Bonnet relations are developed in
\citet{amelunxen2014}.

For a cone $C\subset\mathbb R^3$, let
$C^\circ=\{y:y^{\mathsf T}x\leq0
\ \text{for all }x\in C\}$
denote its polar. Let $\Omega(C)$ be its solid
angle, namely the surface area of
$C\cap\mathbb S^2$.


\begin{proposition}\label{prop:sec5-solid-angle-weights}
For the proper, full-dimensional cone
$\mathcal C_0=J_0^{1/2}K$, the four weights in
\eqref{eq:sec5-chibar} satisfy
\begin{equation}\label{eq:sec5-extreme-weights}
 w_3=\frac{\Omega(\mathcal C_0)}{4\pi},
 \qquad
 w_0=\frac{\Omega(\mathcal C_0^\circ)}{4\pi},
\end{equation}
and
\begin{equation}\label{eq:sec5-middle-weights}
 w_1=\frac12-w_3,
 \qquad
 w_2=\frac12-w_0.
\end{equation}
Thus two solid-angle calculations determine all four weights.
\end{proposition}

The first identities are the extreme conic intrinsic volumes; the
remaining identities follow from the conic Gauss--Bonnet relations in
three dimensions. Appendix~\ref{app:global_lrt} provides the geometric details and the
one-dimensional integral used for their numerical evaluation.

\subsection{Solid-angle calculation and numerical quadrature}
\label{subsec:global_solid_angle}
Use the Lorentz transformation $t=L\gamma$ from
Section~\ref{subsec:lorentz_representation}. The Lorentz cone
$\mathcal L_3$ is self-dual under the Euclidean inner product. Hence
\begin{equation}\label{eq:sec5-angle-maps}
 \begin{aligned}
 \mathcal C_0&=A_+\mathcal L_3,
 &A_+&=J_0^{1/2}L^{-1},\\
 \mathcal C_0^\circ&=A_-\mathcal L_3,
 &A_-&=-J_0^{-1/2}L^{\mathsf T}.
 \end{aligned}
\end{equation}
For either invertible map $A$, a boundary-ray parameterization is
\begin{equation}\label{eq:sec5-boundary-rays}
 \begin{aligned}
 r(\varphi)&=(1,\cos\varphi,\sin\varphi)^{\mathsf T},\\
 &0\leq\varphi\leq2\pi.
 \end{aligned}
\end{equation}
Set $b_A(\varphi)=Ar(\varphi)/\|Ar(\varphi)\|$ and
$a_A=Ae_1/\|Ae_1\|$, where $e_1=(1,0,0)^{\mathsf T}$.
The spherical boundary formula yields
\begin{equation}\label{eq:sec5-angle-integral}
 \begin{aligned}
 \Omega(A\mathcal L_3)
 &=\left|\int_0^{2\pi}q_A(\varphi)\,d\varphi\right|,\\
 q_A(\varphi)
 &=\frac{a_A^{\mathsf T}
 [b_A(\varphi)\times b_A'(\varphi)]}
 {1+a_A^{\mathsf T}b_A(\varphi)}.
 \end{aligned}
\end{equation}
This smooth one-dimensional integral is evaluated for $A_+$ and
$A_-$ using deterministic quadrature. Appendix~\ref{app:global_lrt} derives the formula,
its orientation, and stable equivalent expressions. In computation,
the resulting weights should be checked for nonnegativity, unit sum,
and the two half-sum identities in
\eqref{eq:sec5-middle-weights}.

\subsection{Plug-in weights and asymptotic critical values}
\label{subsec:global_plugin}
Evaluate the analytical information from Section~\ref{sec:information}
at the restricted null estimator $(\widetilde\lambda_n,0)$.
Let $\widehat J_n$ denote the Schur complement of the resulting
normalized information matrix. For each positive-definite $J$,
write $w_j(J)$ for the exact conic intrinsic-volume weight of
$J^{1/2}K$. Let
$\widehat w_n=(\widehat w_{0,n},\ldots,\widehat w_{3,n})$
denote the weights computed by deterministic solid-angle quadrature
using $\widehat J_n$.
For $t>0$, the plug-in upper-tail probability is
\begin{equation}\label{eq:sec5-plugin-p}
 \widehat p_n(t)=
 \sum_{j=1}^{3}\widehat w_{j,n}
 \Pr(\chi_j^2\geq t).
\end{equation}
At $t=0$, the inclusive upper-tail probability is one; the limiting
law has an atom of size $w_0=w_0(J_0)$ at zero. The positive-tail
formula above is not the inclusive upper-tail probability at zero.

\begin{proposition}\label{prop:sec5-plugin-validity}
Suppose the assumptions of Theorem~\ref{thm:sec5-global-lrt}
hold, the restricted-null analytical information estimator is
consistent so that $\widehat J_n\xrightarrow{p}J_0\succ0$, and
$\widehat J_n$ is positive definite with probability tending to one,
and the numerical weights satisfy
\begin{equation}\label{eq:sec5-numerical-error}
 \max_{0\leq j\leq3}
 \bigl|\widehat w_{j,n}-w_j(\widehat J_n)\bigr|
 =o_p(1).
\end{equation}
Then
\[
 \widehat w_{j,n}\xrightarrow{p}w_j(J_0),
 \qquad j=0,1,2,3.
\]
For a fixed $\alpha\in(0,1-w_0(J_0))$, let
$\widehat c_{1-\alpha}$ be the upper $\alpha$ critical value of
the plug-in chi-bar-square law. The test rejecting when
$\mathrm{LRT}_{\mathrm{omni},n}>\widehat c_{1-\alpha}$ has asymptotic size
$\alpha$.
\end{proposition}

The result follows from continuity of the solid angles under
nonsingular linear transformations, the $o_p(1)$ quadrature error
condition, and convergence of the strictly positive mixture
critical value. Appendix~\ref{app:global_lrt} gives the proof.
A fixed numerical integration tolerance alone is not an
asymptotically vanishing error condition. The true-$J_0$ weights
provide an oracle benchmark, whereas the restricted-null plug-in
weights define the feasible procedure. Finite-sample calibration
is examined by simulation.


\subsection{Compact-uniform four-family calibration}
\label{subsec:sec5-uniform-omnibus}
The preceding asymptotic law is a fixed-parameter statement. A
uniform guarantee over a class of omnibus-null nuisance parameters
requires more than pointwise local asymptotic normality. We state the
corresponding result for the explicit model and design classes in the
completed regularity theory. The covariance vertex is fixed under this
null, so the uniformity concerns nuisance variation; uniformity across
covariance ranks arises in the pointwise, contrast, and confidence-region
problems addressed in later sections.

Let $\Lambda_+\subset\mathbb R^{p_\lambda}$ be a fixed compact convex
admissible nuisance fitting set, with a compact true class
$\Lambda_0\subset\operatorname{int}(\Lambda_+)$. Fit over
\begin{equation}\label{eq:sec5-compact-fit}
\Theta_+=\Lambda_+\times
\left\{\gamma\in K:
\operatorname{tr}\{\Gamma(\gamma)\}\leq M\right\},
\qquad M>0,
\end{equation}
so the artificial upper covariance boundary is separated from the
null vertex. The clusters are independent, with bounded sizes and a
finite collection of deterministic design types having limiting
frequencies. The collection contains a positive-frequency informative
design (or the explicitly identifying grouped-design extension), and
the full information is uniformly positive definite. The nuisance
dimension is arbitrary but fixed. For Gaussian models, the residual
variance stays in a compact subset of $(0,\infty)$; for NB2 the
unknown size does likewise; and for compound Poisson--Gamma models
both unknown dispersion and power remain in compact subsets of
$(0,\infty)$ and $(1,2)$, respectively. The specified Gaussian
identity-link, and Poisson, NB2, and compound Poisson--Gamma log-link
models satisfy the required marginal likelihood regularity under
these informative designs.

\begin{corollary}\label{cor:sec5-uniform-omnibus}
Under the preceding compact fitted and informative-design
conditions, suppose the global full and null likelihood suprema are
attained (or approximated to a uniform $o_p(1)$ objective gap),
and let $\widehat J_n$ be a uniformly consistent positive-definite
restricted-null information estimator. Suppose further that the
computed weights satisfy, for each $\varepsilon>0$,
\begin{equation}\label{eq:sec5-uniform-weights}
\sup_{\lambda\in\Lambda_0}
\Pr_{(\lambda,0)}\!\left\{
\max_{0\leq j\leq3}
\left|\widehat w_{j,n}-w_j(\widehat J_n)\right|
>\varepsilon\right\}\longrightarrow0.
\end{equation}
For any fixed $0<\alpha<1/2$, the positive plug-in
chi-bar-square critical value $\widehat c_{1-\alpha}$ obeys
\begin{equation}\label{eq:sec5-uniform-size}
\sup_{\lambda\in\Lambda_0}
\left|
\Pr_{(\lambda,0)}\!\left\{
\mathrm{LRT}_{\mathrm{omni},n}>
\widehat c_{1-\alpha}\right\}-\alpha
\right|\longrightarrow0.
\end{equation}
Thus the omnibus test has compact-uniform asymptotic size
$\alpha$ for the stated four-family classes.
\end{corollary}

\begin{proof}
The verified common likelihood envelopes, identification, and
information bounds give uniform root-$n$ localization, a
true-centered quadratic expansion, and a uniform efficient-score
central limit theorem at the covariance vertex. Because the local
null and alternative covariance sets are respectively $\{0\}$ and
$K$, joint likelihood maximization yields the Gaussian projection
law in \eqref{eq:sec5-lrt-limit} uniformly along every sequence
$\lambda_n\in\Lambda_0$. Continuity of the
information-transformed solid angles and uniform consistency of
$\widehat J_n$ and the numerical weights imply uniform convergence
of the mixture weights. The polar solid-angle probability
$w_0(J)$ is strictly below $1/2$ for every proper three-dimensional
information-transformed PSD cone. Compactness of the information
class makes its separation from $1/2$ uniform. Therefore
$0<\alpha<1/2$ yields critical values uniformly bounded away
from zero, where the limiting mixture distribution is continuous
and strictly increasing. Uniform Gaussian approximation, critical-
value consistency, and this positive-threshold continuity establish
\eqref{eq:sec5-uniform-size}.
\end{proof}

The result concerns the \emph{ideal marginal-likelihood} test with
correctly constrained global suprema. Uniform $o_p(1)$ numerical
errors in the complete maximized log likelihoods and in critical
values are sufficient for transfer to an implemented procedure;
fixed quadrature or optimization tolerances alone do not verify
those requirements. The hypotheses exclude growing cluster sizes,
unidentified designs, nuisance parameters approaching their
physical endpoints, and sequences with singular full Fisher
information. They do \emph{not} fix either compound
Poisson--Gamma dispersion or power.

\subsection{Comparison with conventional chi-square calibration}
\label{subsec:global_chisq_comparison}
A conventional three-degree-of-freedom reference ignores the
positive-semidefinite restriction and treats the three covariance
coordinates as locally unrestricted. At the vertex, however, the
local alternative is the proper cone $K$, not $\mathbb R^3$.
Consequently, the $\chi_3^2$ approximation does not reproduce the
chi-bar-square null law. It nevertheless provides a conservative
asymptotic reference.

\begin{proposition}\label{prop:sec5-chisq-conservative}
Under the assumptions of Theorem~\ref{thm:sec5-global-lrt},
let $T=\|\Pi_{\mathcal C_0}(G)\|^2$ with
$G\sim N_3(0,I_3)$. Then, for every $t\geq0$,
\begin{equation}\label{eq:sec5-chisq-order}
 \Pr(T>t)\leq\Pr(\chi_3^2>t).
\end{equation}
Thus the test rejecting when $\mathrm{LRT}_{\mathrm{omni},n}$ exceeds the
upper $\alpha$ critical value of $\chi_3^2$ has asymptotic
rejection probability at most $\alpha$ under $H_0$.
\end{proposition}

Indeed, Moreau's orthogonal decomposition gives
\[
 \|G\|^2=\|\Pi_{\mathcal C_0}(G)\|^2
 +\|\Pi_{\mathcal C_0^\circ}(G)\|^2,
\]
so $T\leq\|G\|^2$ almost surely, while
$\|G\|^2\sim\chi_3^2$. The complete proof is in
Appendix~\ref{app:global_lrt}.
We compare the oracle and plug-in cone calibrations with this
conservative conventional reference through empirical rejection
rates, power, and representative distributional diagnostics.

\subsection{Relationship to previous variance-component
testing methods}
\label{subsec:global_prior_methods}

Likelihood-ratio inference under boundary constraints
has a substantial theoretical foundation
\citep{chant1974asymptotic,self1987asymptotic,
shapiro1985}. In mixed-effects models,
\citet{stram1994variance} investigated
variance-component testing for longitudinal data,
while \citet{zhang2008variance} considered
variance-component testing in generalized linear
mixed models.

The chi-bar-square structure is an established consequence of
likelihood-ratio theory under convex-cone restrictions
\citep{self1987asymptotic,shapiro1985,amelunxen2014}. Our
specialization uses the random-intercept-and-slope
PSD cone, the nuisance-adjusted information metric,
and the information calculations developed in
Section~\ref{sec:information}.

The solid-angle formulation avoids approximating
the covariance cone by an orthant and provides
a direct numerical route to its exact asymptotic
weights. The geometric derivations and calibration arguments are given in
Appendix~\ref{app:global_lrt}; the comparison with earlier
variance-component methods is confined to the present subsection.
